\section{Limita a spojitost funkce}
\subsection{Základní pojmy}
\begin{definice}
\textbf{Funkce} $f$ \textbf{jedné reálné proměnné} (dále jen \textbf{funkce}) je zobrazení
$f: M\to \mathbb R $, kde $M$ je podmnožinou množiny reálných čísel.
\end{definice}

\begin{definice}
	Funkce $f: J\to \mathbb R $ je \textbf{rostoucí} na intervalu $J$, jestliže pro
	každou dvojici $x_1, x_2 \in J$, $x_1 < x_2$ platí nerovnost $f(x_1)<f(x_2)$. Analogicky
	definujeme funkci \textbf{klesající} (\textbf{neklesající}, \textbf{nerostoucí}) na intervalu $J$.
\end{definice}

\begin{definice}
\textbf{Monotónní funkcí} (resp. \textbf{ryze monotónní funkcí}) na intervalu $J$ rozumíme funkci,
která je neklesající nebo nerostoucí (resp. rostoucí nebo klesající) na $J$.
\end{definice}

\begin{definice}
Nechť $f$ je funkce a $M \subset \mathcal D(f)$. Řekneme, že funkce $f$ je
\begin{itemize}
  \item \textbf{lichá}, jestliže pro každé $x \in \mathcal D(f)$ platí $-x \in \mathcal D(f)$ a
  $f(-x) = -f(x)$,
  \item \textbf{sudá}, jestliže pro každé $x \in \mathcal D(f)$ platí $-x \in \mathcal D(f)$ a
  $f(-x) = f(x)$,
  \item \textbf{periodická s periodou} $a \in \mathbb{R}$, $a > 0$, jestliže pro každé
  $x \in \mathcal D(f)$ platí $x + a \in \mathcal D(f)$, $x - a \in \mathcal D(f)$ a
  $f(x + a) = f(x - a) = f(x)$,
  \item \textbf{shora omezená} na $M$, jestliže existuje číslo $K \in \mathbb{R}$ takové,
  že pro všechna $x \in M$ je $f(x) \le K$,
  \item \textbf{zdola omezená} na $M$, jestliže existuje číslo $K \in \mathbb{R}$ takové,
  že pro všechna $x \in M$ je $f(x) \ge K$,
  \item \textbf{omezená} na $M$, jestliže existuje číslo $K \in \mathbb{R}$ takové,
  že pro všechna $x \in M$ je $|f(x)| \le K$,
  \item \textbf{konstantní} na $M$, jestliže pro všechna $x, y \in M$ platí
  $f(x) = f(y)$.
\end{itemize}
\end{definice}

\subsection{Limita funkce}

\begin{definice}
	Nechť $c \in \mathbb R $. Potom \textbf{okolím bodu} $c$ rozumíme každou množinu tvaru
	$B(c, \varepsilon) = (c-\varepsilon, c+\varepsilon)$, kde $\varepsilon \in \mathbb R , \varepsilon > 0$.
	\textbf{Okolím} bodu $\infty$ rozumíme každou množinu tvaru $B(\infty, \varepsilon) = \left( \frac{1}{\varepsilon}, \infty \right ) $,
	kde $\varepsilon \in \mathbb R , \varepsilon > 0$. \textbf{Okolím} bodu $-\infty$
	rozumíme každou množinu $B(-\infty, \varepsilon)= \left( -\infty, -\frac{1}{\varepsilon} \right ) $, kde $\varepsilon \in \mathbb R $, $\varepsilon >0$.
\end{definice}

\begin{definice}
	Nechť $c \in \mathbb R $. Potom \textbf{prstencovým okolím} bodu $c$ rozumíme
	každou množinu tvaru $(c-\varepsilon, c+\varepsilon) \smallsetminus \left \{ c \right \} $,
	kde $\varepsilon \in \mathbb R ,\varepsilon > 0$, a kterou značíme
	$P(c,\varepsilon)$. \textbf{Prstencovým okolím} bodu $\infty$ rozumíme každou množinu tvaru $\left( \frac{1}{\varepsilon},\infty \right ) $,
	kde $\varepsilon \in \mathbb R , \varepsilon > 0$, a kterou značíme $P(\infty, \varepsilon)$.
	\textbf{Prstencovým okolím} bodu $-\infty$ rozumíme každou množinu tvaru
	$\left( -\infty, -\frac{1}{\varepsilon} \right ) $, kde $\varepsilon \in \mathbb R ,\varepsilon > 0$,
	a kterou značíme $P(-\infty, \varepsilon)$.
\end{definice}

\begin{pozorovani}
	Nechť $c_1,c_2 \in \mathbb R ^*, c_1 \ne c_2$. Potom existuje $\varepsilon \in \mathbb R , \varepsilon > 0$
	takové, že $B(c_1, \varepsilon) \cap B(c_2,\varepsilon) = \emptyset$.
\end{pozorovani}

\begin{definice}
	Řekneme, že číslo $A \in \mathbb R ^*$ je \textbf{limitou funkce $f$ v bodě $c \in \mathbb R ^*$},
	jestliže
	$$
	    \forall \varepsilon \in \mathbb R, \varepsilon > 0 \, \exists \delta \in \mathbb R , \delta > 0 \, \forall x \in P(c, \delta): f(x)\in B(A,\varepsilon).
	$$
\end{definice}

\begin{veta}\label{3.1}% 3.1
	Nechť funkce $f$ má v bodě $c\in \mathbb R ^*$ limitu $A\in \mathbb R ^*$ a limitu
	$B\in \mathbb R ^*$. Potom platí $A=B$.
\end{veta}

\begin{proof}
	Pro spor předpokládejme, že $A\ne B$. Nalezneme $\varepsilon \in \mathbb R , \varepsilon >0$
	takové, že $B(A,\varepsilon)\cap B(B,\varepsilon) = \emptyset$. Nalezneme $\delta_1, \delta_2 \in \mathbb R ,
	\delta_1, \delta_2 > 0$ takové, že platí
	$$
	    \forall x \in P(c,\delta_1): f(x) \in B(A,\varepsilon),\\
		\forall x \in P(c,\delta_2): f(x) \in B(B,\varepsilon).
	$$
	Položme $\delta_3 = \min \left \{ \delta_1, \delta_2 \right \} $. Zvolme $z\in P(c,\delta_3)$.
	Potom platí $c\in P(c,\delta_1)$, a tedy $f(z)\in B(A,\varepsilon).$ Podobně
	máme $z\in P(c,\delta_2)$, a tedy $f(z)\in B(B,\varepsilon).$ Odtud plyne
	$$
	    f(z)\in B(A, \varepsilon) \cap B(B,\varepsilon) = \emptyset,
	$$
	což je spor.
\end{proof}

\begin{poznamka}\,
\begin{enumerate}[(1)]
	\item Pokud existuje limita funkce $f$ v bodě $c\in \mathbb R ^*$, pak ji značíme symbolem
	$\lim_{x\to c} f(x)$.
	\item Pokud existuje $\lim_{x\to c} f(x)$, potom je $f$ definována na jistém prstencovém okolí
	bodu $c$.
	\item Limitu funkce můžeme počítat ve vlastním bodě $c$, tj. $c\in \mathbb R $, nebo v nevlastním
	bodě $c$, tj. $c\in \left \{ \infty, -\infty \right \} $. Výsledek může být opět vlastní nebo
	nevlastní.
	\item Nechť $c, A\in \mathbb R .$ Potom platí, že $\lim_{x\to c}f(x)=A$ právě tehdy, když
	$$
	    \forall \varepsilon \in \mathbb R , \varepsilon > 0 \, \exists \delta \in \mathbb R , \delta > 0 \, \forall x \in \mathbb R , 0 < |x-c|<\delta:|f(x)-A|<\varepsilon.
	$$
	\item Nechť $c,A\in \mathbb R , K\in \mathbb R , K>0$. Potom $\lim_{x\to c}f(x)=A$ právě tehdy, když platí
	$$
	\forall \varepsilon \in \mathbb R , \varepsilon > 0 \, \exists \delta \in \mathbb R , \delta > 0 \, \forall x \in \mathbb R , 0 < |x-c|<\delta:|f(x)-A|<K\varepsilon.
	$$
\end{enumerate}
\end{poznamka}

\begin{definice}
	Nechť $c\in \mathbb R $.
	\begin{itemize}
  \item \textbf{Pravým okolím} bodu $c$ rozumíme každý interval
  $[c, c + \varepsilon)$, kde $\varepsilon > 0$, a který značíme $B^+(c,\varepsilon)$,
  \item \textbf{levým okolím} bodu $c$ rozumíme každý interval
  $(c - \varepsilon, c]$, kde $\varepsilon > 0$, a který značíme $B^-(c,\varepsilon)$,
  \item \textbf{pravým prstencovým okolím} bodu $c$ rozumíme každý interval
  $(c, c + \varepsilon)$, kde $\varepsilon > 0$, a který značíme $P^+(c,\varepsilon)$,
  \item \textbf{levým prstencovým okolím} bodu $c$ rozumíme každý interval
  $(c - \varepsilon, c)$, kde $\varepsilon > 0$, a který značíme $P^-(c,\varepsilon)$.
\end{itemize}
Pokračujme s definicí pro nevlastní hodnoty.
\begin{itemize}
  \item \textbf{Pravým okolím bodu $-\infty$} rozumíme každý interval
  $\left(-\infty, -\frac{1}{\varepsilon}\right)$, kde $\varepsilon > 0$,
  a který značíme $B^+(-\infty,\varepsilon)$,
  \item \textbf{levým okolím bodu $\infty$} rozumíme každý interval
  $\left(\frac{1}{\varepsilon}, \infty\right)$, kde $\varepsilon > 0$,
  a který značíme $B^-(\infty,\varepsilon)$,
  \item \textbf{pravým prstencovým okolím bodu $-\infty$} rozumíme každý interval
  $\left(-\infty, -\frac{1}{\varepsilon}\right)$, kde $\varepsilon > 0$,
  a který značíme $P^+(-\infty,\varepsilon)$,
  \item \textbf{levým prstencovým okolím bodu $\infty$} rozumíme každý interval
  $\left(\frac{1}{\varepsilon}, \infty\right)$, kde $\varepsilon > 0$,
  a který značíme $P^-(\infty,\varepsilon)$.
\end{itemize}
\end{definice}

\begin{definice}
Nechť $A \in \mathbb{R}^\star$, $c \in \mathbb{R} \cup \{-\infty\}$.
Řekneme, že funkce $f$ má v bodě $c$ \textbf{limitu zprava} rovnou $A$, jestliže
\[
\forall \varepsilon \in \mathbb{R},\ \varepsilon > 0\ \exists \delta \in \mathbb{R},\
\delta > 0\ \forall x \in P^+(c,\delta) : f(x) \in B(A,\varepsilon).
\]
Analogicky definujeme pojem \textbf{limity zleva} v bodě
$c \in \mathbb{R} \cup \{\infty\}$.
\end{definice}

\begin{poznamka}
	Věta \ref{3.1} platí analogicky i pro limitu zleva a limitu zprava. Značení: $\lim_{x\to c_-}f(x)$, resp.
	$\lim_{x\to c_+}f(x)$.
\end{poznamka}

\begin{priklad}\,
\begin{itemize}
	\item Nechť $c\in \mathbb R ^*, A \in \mathbb R ,\delta_0\in \mathbb R ,\delta_0>0$ splňují
	$\forall x \in P(c,\delta_0):f(x)=A.$
	Potom $\lim_{x\to c}f(x)=A.$
	\item Nechť $c\in \mathbb R ^*$. Potom platí $\lim_{x\to c}x=c$. Zvolme
	$\varepsilon \in \mathbb R ,\varepsilon >0$. Položme $\delta=\varepsilon$. Potom platí
	$P(c,\delta)=P(c,\varepsilon)\subseteq B(c,\varepsilon),$ a tedy
	$\forall x \in P(c,\delta):x\in B(c,\varepsilon)$.
\end{itemize}
\end{priklad}

\begin{pozorovani}
	Nechť $f,g$ jsou funkce, $c\in \mathbb R ^*, \delta_0 \in \mathbb R , \delta_0 > 0$ a platí
	$$
	    \forall x\in P(c,\delta_0): f(x)=g(x).
	$$
	Potom $\lim_{x\to c}f(x)=A\in \mathbb R ^*$ právě tehdy, když $\lim_{x\to c}g(x)=A\in \mathbb R ^*$.
\end{pozorovani}

\begin{poznamka}
	Limita funkce v bodě má \uv{lokální charakter}.
\end{poznamka}

\begin{definice}
	Řekneme, že funkce $f$ je \textbf{spojitá v bodě} $c\in \mathbb R $, jestliže platí
	$\lim_{x\to c} f(x)=f(c)$.
\end{definice}

\begin{poznamka}
	Funkce je spojitá v bodě $c\in \mathbb R $ právě tehdy, když platí
	$$
	    \forall \varepsilon \in \mathbb R , \varepsilon > 0 \,\exists \delta \in \mathbb R , \delta > 0 \, \forall x \in B(c,\delta):f(x)\in B(f(c),\varepsilon).
	$$
\end{poznamka}

\begin{definice}
Řekneme, že funkce $f$ je \textbf{spojitá zprava} (resp. \textbf{zleva}) v bodě $c \in \mathbb{R}$,
jestliže platí
\[
\lim_{x \to c^+} f(x) = f(c)
\quad
(\text{resp. } \lim_{x \to c^-} f(x) = f(c)).
\]
\end{definice}

\begin{definice}
Nechť $J \subset \mathbb{R}$ je interval. Funkce $f \colon J \to \mathbb{R}$
je \textbf{spojitá na intervalu} $J$, jestliže platí:
\begin{itemize}
  \item $f$ je spojitá zprava v každém bodě $J$, který není pravým koncovým bodem $J$,
  \item $f$ je spojitá zleva v každém bodě $J$, který není levým koncovým bodem $J$.
\end{itemize}
\end{definice}

\begin{poznamka}
	Nechť $c\in \mathbb R , A\in \mathbb R ^*$. Potom platí, že $\lim_{x\to c}f(x)=A$ právě tehdy,
	když $\lim_{x\to c_-}f(x)=A$ a $\lim_{x\to c_+}f(x)=A$.
	\par Implikace $\implies$ je zřejmá. Implikace $\impliedby$: Zvolme $\varepsilon \in \mathbb R ,\varepsilon > 0$.
	Nalezneme $\delta _1\in \mathbb R , \delta_1 >0$ takové, že
	$$
	    \forall x\in P_1(c,\delta_1) : f(x)\in B(A,\varepsilon).
	$$
	Podobně nalezneme $\delta_2\in \mathbb R , \delta_2 >0$ takové, že
	$$
	    \forall x\in P_2(c,\delta_2):f(x)\in B(A,\varepsilon).
	$$
	Položme $\delta_3=\min \left \{ \delta_1, \delta_2 \right \} $. Potom máme
	$$
	    P(c,\delta_3) \subseteq P_+(c,\delta_1) \cup P_-(c,\delta_2).
	$$
	Pro každé $x \in P(c,\delta_3)$ platí $f(x)\in B(A,\varepsilon)$.
\end{poznamka}

\subsection{Věty o limitách}
\begin{veta}\label{veta3.2}
	Nechť funkce $f$ má v bodě $c\in \mathbb R ^*$ vlastní limitu. Pak existuje
	$\delta\in \mathbb R ,\delta >0$ takové, že $f$ je na $P(c,\delta)$ omezená.
\end{veta}

\begin{proof}
	Označme $A=\lim_{x\to c}f(x)$. Potom platí $A\in \mathbb R $. Položme $\varepsilon =1$. K němu
	nalezneme $\delta\in \mathbb R ,\delta > 0$ takové, že pro každé $x\in P(c,\delta): f(x)\in B(A,\varepsilon)$,
	neboli
	$$
	    \forall x \in P(c,\delta):f(x)\in(A-1, A+1).
	$$
	Tedy $f$ je omezená na $P(c,\delta)$.
\end{proof}

\begin{veta}[Aritmetika limit]\label{veta3.3} % věta 3.3
Nechť $A, B, c \in \mathbb{R}^\star$, $\lim_{x \to c} f(x) = A$ a
$\lim_{x \to c} g(x) = B$. Potom platí:
\begin{itemize}
  \item[(a)]
  $
  \lim_{x \to c} \bigl(f(x) + g(x)\bigr) = A + B,
  $
  pokud je výraz na pravé straně definován,
  \item[(b)]
  $
  \lim_{x \to c} f(x) g(x) = AB,
  $
  pokud je výraz na pravé straně definován,
  \item[(c)]
  $
  \lim_{x \to c} \frac{f(x)}{g(x)} = \frac{A}{B},
  $
  pokud je výraz na pravé straně definován.
\end{itemize}
\end{veta}

\begin{proof}\,
\begin{enumerate}
	\item[(b)] Případ $A,B \in \mathbb R $: Podle věty \ref{veta3.2}
	nalezneme $K\in \mathbb R , K > 0$ takové, že existuje $\delta_0 \in \mathbb R , \delta_0 > 0$ splňující
	$$
	    \forall x \in P(c,\delta_0):|f(x)| < K.
	$$
	Zvolme $\varepsilon \in \mathbb R, \varepsilon > 0$. Nalezneme $\delta_1,\delta_2 \in \mathbb R ,\delta_1, \delta_2 > 0$
	taková, že platí
	\begin{align*}
		&\forall x \in P(c,\delta_1): |f(x)-A|<\varepsilon, \\
		&\forall x \in P(c,\delta_2): |g(x)-B|< \varepsilon.
	\end{align*}
	Položme $\delta = \min \left \{ \delta_0, \delta_1, \delta_2 \right \} $.Pro každé $x\in P(c,\delta)$ platí
	\begin{align*}
	|f(x)g(x)-AB| & = |f(x)g(x)-f(x)B+f(x)B-AB| \\
	& \leq |f(x)(g(x)-B)| + |(f(x)-A)B|++
	& < K\varepsilon + \varepsilon |B| = (K+|B|)\varepsilon.
\end{align*}
Odtud plyne $\lim_{x\to c} f(x)g(x)=AB.$
\item[(c)] Případ $A = \infty, B\in \mathbb R , B> 0$. Potom $\frac{A}{B}\infty$. Zvolme $\varepsilon \in \mathbb R ,\varepsilon > 0$.
K němu nalezneme $\delta_1, \delta_2 \in \mathbb R , \delta_1, \delta_2 > 0$ taková, že platí
\begin{align*}
\forall x \in P(c,\delta_1): f(x)\in B(\infty, \frac{1}{2B}\varepsilon),\\
\forall x \in P(c,\delta_2): g(x) \in B(B,\frac{1}{2}B).
\end{align*}
    Položme $\delta = \min \left \{ \delta_1, \delta_2 \right \} $. Pro každé $x \in P(c,\delta)$ platí
    $$
    \frac{f(x)}{g(x)}> \frac{\frac{2B}{\varepsilon}}{\frac{3}{2}B} = \frac{4}{3\varepsilon} > \frac{1}{\varepsilon},
    $$
    tj. $\frac{f(x)}{g(x)} \in B(\infty, \varepsilon).$ \qedhere
\end{enumerate}
\end{proof}

\begin{poznamka}
	Věta \ref{veta3.3} platí analogicky i pro jednostranné limity.
\end{poznamka}

\begin{dusledek}\label{duslvety3.3}
	Nechť $c\in \mathbb R $ a $f,g$ jsou spojité funkce v bodě $c$. Potom také $f+g$ a $fg$ jsou spojité
	v bodě $c$. Pokud navíc $g(c)\ne 0$, pak i funkce $\frac{f}{g}$ je spojitá v $c$.
\end{dusledek}

\begin{priklad}
	Funkce $x\mapsto c, x \in \mathbb R $ a $x\mapsto x, x \in \mathbb R $ jsou spojité na $\mathbb R $.
	Odtud a z důsledku \ref{duslvety3.3} plyne, že každý polynom je spojitý v každém bodě z $\mathbb R $.
\end{priklad}

\begin{veta}\label{veta3.4} % věta 3.4
Nechť $A \in \mathbb{R}^\star$, $A > 0$, $c \in \mathbb{R}^\star$,
$\lim_{x \to c} g(x) = 0$, $\lim_{x \to c} f(x) = A$ a existuje
$\eta \in \mathbb{R}$, $\eta > 0$, takové, že pro každé
$x \in P(c,\eta)$ platí $g(x) > 0$.
Potom platí
\[
\lim_{x \to c} \frac{f(x)}{g(x)} = \infty.
\]
\end{veta}

\begin{proof}
	Rozlišíme dva případy
	\begin{enumerate}[$i$.]
		\item $A\in \mathbb R $: Zvolme $\varepsilon \in \mathbb R , \varepsilon > 0$. K němu
		nalezneme $\delta_1, \delta_2 \in \mathbb R , \delta_1, \delta_2 > 0$ taková, že
		\begin{align*}
		&\forall x \in P(c, \delta_1): |f(x)-A|<\frac{1}{2}A,\\
		&\forall x \in P(c,\delta_2): |g(x) - 0| = |g(x)| < \frac{1}{2} A\varepsilon.
		\end{align*}
		Položme $\delta = \min \left \{ \delta_1, \delta_2, \eta \right \}. $ Potom pro každé $x \in P(c,\delta)$
		platí
		$$
  \frac{f(x)}{g(x)} > \frac{\frac{1}{2}A}{\frac{1}{2}A\varepsilon}=\frac{1}{\varepsilon},
		$$
		tj. $\frac{f(x)}{g(x)} \in B(\infty, \varepsilon)$. Potom tedy $\lim_{x\to c} \frac{f(x)}{g(x)} = \infty$.
		\item $A = \infty$: Zvolme $\varepsilon \in \mathbb R , \varepsilon > 0$. K němu
		nalezneme $\delta_1, \delta_2\in \mathbb R , \delta_1, \delta_2 >0$ taková, že
		\begin{align*}
			&\forall x \in P(c,\delta_1): f(x) \in B(\infty, \varepsilon), \\
			&\forall x \in P(c,\delta_2): G(x) \in B(0, 1).
        \end{align*}
        Položme $\delta = \min \left \{ \delta_1, \delta_2,\eta \right \} $. Potom pro každé $x \in P(c,\delta)$
        platí
        $$
        \frac{f(x)}{g(x)}> \frac{\frac{1}{\varepsilon}}{1} = \frac{1}{\varepsilon}.
        $$
        Odtud plyne $\lim_{x\to c} \frac{f(x)}{g(x)}=\infty$. \qedhere
	\end{enumerate}
\end{proof}

\begin{veta}[Limita a uspořádání]\label{veta3.5} % 3.5
Nechť $c \in \mathbb{R}^\star$ a funkce $f$ a $g$ mají limity v $c$.
\begin{itemize}
  \item[(a)] Nechť
  \[
  \lim_{x \to c} f(x) > \lim_{x \to c} g(x).
  \]
  Pak existuje $\delta \in \mathbb{R}$, $\delta > 0$, takové, že platí
  \[
  \forall x \in P(c,\delta) : f(x) > g(x).
  \]
  \item[(b)] Nechť existuje $\delta \in \mathbb{R}$, $\delta > 0$, takové, že platí
  \[
  \forall x \in P(c,\delta) : f(x) \le g(x).
  \]
  Potom platí
  \[
  \lim_{x \to c} f(x) \le \lim_{x \to c} g(x).
  \]
\end{itemize}
\end{veta}

\begin{proof}
	Označme $A = \lim_{x\to c} f(x), B = \lim_{x\to c} g(x)$.
	\begin{enumerate}[(a)]
		\item Nalezneme $\varepsilon \in \mathbb R , \varepsilon > 0$
		takové, že $B(A,\varepsilon) \cap B(B,\varepsilon) = \emptyset.$ Potom pro každé $y_1 \in B(A,\varepsilon)$
		a $y_2\in B(B,\varepsilon)$ platí $y_1 > y_2$. Nalezneme $\delta_1, \delta_2 \in \mathbb R , \delta_1, \delta_2 > 0$
		taková, že
		\begin{align*}
			& \forall x \in P(c,\delta_1):f(x)\in B(A,\varepsilon), \\
			& \forall x \in P(c,\delta_2):g(x)\in B(B,\varepsilon).
\end{align*}
    Položme $\delta = \min \left \{ \delta_1, \delta_2 \right \} $. Potom pro každé $x \in P(c,\delta)$
    platí $f(x)\in B(A,\varepsilon),g(x)\in B(B,\varepsilon)$, a tedy $f(x) > g(x)$.
    \item Pro spor předpokládejme, že $A > B$. Podle již dokázaného bodu $(a)$ nalezneme
    $\eta \in \mathbb R , \eta > 0$ takové, že
    $$
        \forall x \in P(c,\eta): f(x) > g(x).
    $$
    Zvolme $z \in P(c, \min \left \{ \delta, \eta \right \} )$. Potom platí
    $$
        f(z) > g(z) \geq f(z),
    $$
    což je spor. (Druhá nerovnost z předpokladu.) \qedhere
\end{enumerate}
\end{proof}

\begin{veta}[O dvou strážnících] % 3.6
Nechť $A \in \mathbb{R}^\star$, $\eta \in \mathbb{R}$, $\eta > 0$, a funkce
$f$, $g$ a $h$ splňují:
\begin{itemize}
  \item $\forall x \in P(c,\eta) : f(x) \le h(x) \le g(x)$,
  \item $\lim_{x \to c} f(x) = \lim_{x \to c} g(x) = A$.
\end{itemize}
Potom existuje $\lim_{x \to c} h(x)$ a rovná se $A$.
\end{veta}

\begin{proof}
	Zvolme $\varepsilon \in \mathbb R , \varepsilon > 0$. K němu nalezneme $\delta_1, \delta_2 \in \mathbb R , \delta_1, \delta_2 > 0$
	taková, že
	\begin{align*}
		& \forall x \in P(c,\delta_1):f(x)\in B(A, \varepsilon), \\
		& \forall x \in P(c,\delta_2):g(x)\in B(A, \varepsilon).
\end{align*}
    Položme $\delta = \min \left \{ \delta_1, \delta_2,\eta \right \} $. Potom pro každé
    $x \in B(c,\delta)$ platí $f(x), g(x)\in B(A, \varepsilon)$ a navíc $f(x)\leq h(x)\leq g(x)$,
    tedy $h(x) \in B(A, \varepsilon)$. Odtud plyne, že $\lim_{x\to c} h(x) = A.$
\end{proof}

\begin{poznamka}
	Pokud $A = \infty$, stačí strážník $f$. Podobně pro $A = \infty$.
\end{poznamka}

\begin{dusledek}
	Nechť $c \in \mathbb R ^*, \lim _{x\to c} f(x) = 0$ a existuje $\eta \in \mathbb R ,\eta > 0$
	takové, že funkce $g$ je na $P(c,\eta)$ omezená. Potom $\lim_{x\to c} f(x)g(x)=0.$
\end{dusledek}

\begin{proof}
	Nalezneme $K\in \mathbb R , K > 0$ takové, že $\forall x \in P(c,\eta): |g(x)| <K$.
	Platí také $\lim_{x\to c} |f(x)|=0.$ Potom $\forall x \in P(c, \eta)$ platí
	$$
	    -K |f(x)| \leq f(x)g(x) \leq K|f(x)|.
	$$
	Odtud podle věty o dvou strážnících máme $\lim_{x\to c} f(x)g(x) = 0.$
\end{proof}

\begin{veta}[Limita složené funkce] \label{veta3.7}
Nechť $x_0,y_0\in\R^*$, $f,g$ jsou funkce, $\lim_{x\to x_0} g(x)=y_0$, $\lim_{y\to y_0} f(y)$ existuje a je splněna alespoň jedna z podmínek:
\begin{enumerate}
	\item[(P)] $\exists\eta \in \mathbb R , \eta >0\ \forall x\in P(x_0,\eta): g(x)\ne y_0$,
	\item[(S)] $f$ je spojitá v bodě $y_0$.
\end{enumerate}
Potom platí $$\lim_{x\to x_0} (f\circ g)(x)=\lim_{y\to y_0} f(y).$$
\end{veta}

\begin{proof}
	Označme $A=\lim_{y\to \_0}f(y)$. Chceme $\lim_{x\to x_0} f(g(x))=A$.
	\begin{enumerate}
		\item[(P)] Zvolme $\varepsilon \in \mathbb R , \varepsilon > 0$. K nim nalezneme
		$\psi \in \mathbb R , \psi > 0$ takové, že
		$$\forall y \in P(y_0, \psi): f(y)\in B(A, \varepsilon).$$
		K $\psi$ nalezneme $\delta^\prime \in \mathbb R , \delta ^\prime > 0$ takové, že platí
		$$
		    \forall x \in P(x_0,\delta^\prime): g(x)\in B(y_0,\psi).
		$$
		Nechť $\eta \in \mathbb R,\eta>0 $ splňuje podmínku (P). Položme
		$\delta = \min \left \{ \delta^\prime, \eta \right \} $.
		Pro každé $x \in P(x_0,\delta)$ platí $g(x\in B(y_0,\psi)$ a $g(x)\ne y_0$, tedy $g(x)\in P(y_0, \psi)$.
		Odtud plyne $f(g(x))\in B(A,\varepsilon).$
		\item[(S)] Zvolme $\varepsilon \in \mathbb R , \varepsilon > 0$. K němu díky (S) nalezneme
		$\psi \in \mathbb R , \psi > 0 $ takové, že
		$$
		    \forall y \in B(y_0,\psi) : f(y) \in B(f(y_0), \varepsilon),
		$$
		neboť $A=f(y_0)$. K $\psi$ nalezneme $\delta \in \mathbb R ,\delta > 0$ takové, že
		$$
		    \forall x \in P(x_0, \delta):g(x) \in B(y_0, \psi).
		$$
		Potom pro každé $x \in P(x_0, \delta)$ platí $f(g(x)) \in B(f(y_0),\varepsilon)=B(A,\varepsilon)$. \qedhere
	\end{enumerate}
\end{proof}

\begin{dusledek}
	Jestliže je funkce $g$ spojitá v bodě $c\in \mathbb R $ a funkce $f$ je spojitá v bodě $g(c)$, potom
	je funkce $f\circ g$ spojitá v bodě $c$.
\end{dusledek}

\begin{poznamka}
	Platí i varianty věty \ref{veta3.7} pro jednostranné limity, např.: Nechť $x_0,y_0 \in \mathbb R \smallsetminus \left \{ -\infty \right \} $,
	$f,g$ jsou funkce, $\lim_{x\to x_0^+} g(x)=y_0,\lim_{y\to y_0^+} f(y)$ existuje a je splněna alespoň jedna z následujících
	podmínke:
	\begin{enumerate}
		\item[(P)] $\exists \eta \in \mathbb R , \eta > 0\, \forall x \in P_+(x_0,\eta):g(x)>y_0$,
		\item[(S)] $f$ je spojitá v bodě $y_0$,
\end{enumerate}
    potom $\lim_{x\to x_0^+} f(g(x))=\lim_{y\to y_0^+}f(y)$.
\end{poznamka}

\begin{veta}[Heine]\label{Heinehoveta}
Nechť $c,A\in \mathbb R ^*$ a $f$ je funkce. Pak jsou následující výroky ekvivalentní:
\begin{enumerate}[(i)]
  \item $\lim_{x\to c} f(x)=A$.
  \item Pro každou posloupnost $\{x_n\}$ splňující
  \begin{itemize}
	\item $\forall n\in \mathbb N :x_n\ne c$ a
	\item $\lim_{n\to\infty} x_n=c$
  \end{itemize}
  platí $\lim_{n\to\infty} f(x_n)=A$.
\end{enumerate}
\end{veta}

\begin{proof}\,
\begin{enumerate}
	\item[$\implies$] Nechť $\left \{ x_n \right \} $ je posloupnost splňující
	\begin{itemize}
		\item $\forall n\in \mathbb N : x_n \ne c,$
		\item $\lim_{n \to \infty} x_n = c$.
	\end{itemize}
	Chceme ukázat, že $\lim_{n \to \infty} f(x_n)=A$. Zvolme $\varepsilon \in \mathbb R , \varepsilon > 0$.
	Díky ($i$) nalezneme k $\varepsilon$ číslo $\delta\in \mathbb R , \delta > 0$ takové, že
	$$
	    \forall x \in P(c,\delta):  f(x) \in B(A,\varepsilon).
	$$
	K $\delta$ nalezneme $n_0 \in \mathbb N $ takové, že
	$$
	    \forall n \in \mathbb N , n \geq n_0: x_n \in B(c,\delta),
	$$
	neboť $\lim_{n \to \infty} x_n = c$. Pak máme
	$$
	    \forall n \in \mathbb N , n \geq n_0: x_n \in P(c,\delta).
	$$
	Potom z předchozího plyne
	$$
	    \forall n \in \mathbb N , n \geq n_0: f(x_n) \in B(A,\varepsilon).
	$$
	Odtud $\lim_{n \to \infty} f(x_n)=A.$
	\item[$\impliedby$] Obměnou. Předpokládejme negaci ($i$):
	$$
	    \exists \varepsilon \in \mathbb R , \varepsilon > 0 \, \forall \delta \in \mathbb R , \delta > 0 \, \exists x \in P(c,\delta): \neg(f(x)\in B(A,\varepsilon)).
	$$
	Nechť $\varepsilon \in \mathbb R , \varepsilon > 0$ a podle $\neg$ ($i$) splňuje
	$$
	    \forall \delta \in \mathbb R , \delta > 0 \, \exists x \in P(c,\delta): \neg (f(x)\in B(A,\varepsilon).
	$$
	Pro každé $n \in \mathbb N $ nalezneme $x_n \in P \left( c, \frac{1}{n} \right ) $ takové, že
	$\neg(f(x_n)\in B(A,\varepsilon)$. Potom pro každé $n \in \mathbb N $ platí $x_n \ne c$.
	Dále máme $\lim_{n \to \infty} x_n = c$. Zvolme $\tau \in \mathbb R , \tau > 0$. K němu nalezneme
	$n_0 \in \mathbb N $ splňující $\frac{1}{n_0} < \tau$. Potom pro každé $n\in \mathbb N , n \geq n_0$
	platí
	$$
	    x_n \in P(c, \frac{1}{n}) \subseteq P(c,\frac{1}{n_0})\subseteq B(c,\tau).
	$$
	Neplatí však $\lim_{n \to \infty} f(x_n)=A$, neboť pro každé $n \in \mathbb N $ máme
	$\neq (f(x_n)\in B(A,\varepsilon))$.\qedhere
\end{enumerate}
\end{proof}

\begin{poznamka}
	Lze formulovat varianty Heineovy věty, např.: Nechť $c\in \mathbb R $ a $f$ je funkce.
	Pak jsou výroky ($i$) a ($ii$) ekvivalentní:
	\begin{enumerate}[($i$)]
		\item $f$ je spojitá zprava v bodě $c$,
		\item pro každou posloupnost $\left \{ x_n \right \} $ splňující
		\begin{itemize}
			\item $\forall n \in \mathbb N : x_n \geq c$,
			\item $\lim_{n \to \infty} x_n = c$,
		\end{itemize}
		platí $\lim_{n \to \infty} f(x_n)=f(c)$.
	\end{enumerate}
\end{poznamka}

\begin{priklad}
	Uvažujme funkci $f(x)=\sin \frac{1}{x}, x \in \mathbb R \smallsetminus \left \{ 0 \right \} $.
	Položme $x_n=\frac{1}{n\pi},n\in \mathbb N $, $y_n = \frac{1}{2n\pi + \frac{\pi}{2}}, n \in \mathbb N $.
	Platí $x_n \ne 0, y_n \ne 0$ pro všechna $n \in \mathbb N $. Navíc je
	\begin{align*}
		\lim_{n \to \infty} f(x_n)=\lim_{n \to \infty} \sin(n\pi)=0, \\
		\lim_{n \to \infty} f(y_n) = \lim_{n \to \infty} \left( 2n\pi + \frac{\pi}{2} \right ) =1,
\end{align*}
tedy $\lim_{x \to 0} f(x)$ neexistuje.
\end{priklad}

\begin{veta}[Limita monotónní funkce]\label{veta3.9}
Nechť $a,b\in \mathbb R ^*$, $a<b$, $f$ je monotónní funkce na intervalu $(a,b)$. Potom existují $\lim_{x\to a^+} f(x)$ a $\lim_{x\to b^-} f(x)$, přičemž platí:
\begin{enumerate}[(a)]
  \item je-li $f$ na $(a,b)$ neklesající, pak
  \begin{align*}
	\lim_{x\to a^+} f(x)=\inf f((a,b)) && \textrm{a} & & \lim_{x\to b^-} f(x)=\sup f((a,b)),
  \end{align*}
  \item je-li $f$ na $(a,b)$ nerostoucí, pak
  \begin{align*}
	\lim_{x\to a^+} f(x)=\sup f((a,b)) && \textrm{a} & & \displaystyle \lim_{x\to b^-} f(x)=\inf f((a,b)).
  \end{align*}
\end{enumerate}
\end{veta}

\begin{proof}\,
\begin{enumerate}[(a)]
	\item Označme $m=\inf f((a,b))$. Předpokládejme, že $m\in \mathbb R $. Zvolme $\varepsilon \in \mathbb R , \varepsilon > 0$.
	Podle definice infima není číslo $m+\varepsilon$ dolní závorou množiny $f((a,b))$. Nalezneme $x^\prime \in (a,b)$
	takové, že $m\leq f(x^\prime)< m + \varepsilon$. Nalezneme $\delta \in \mathbb R , \delta  > 0$ takové, že
	$P_+(a,\delta)\subset (a,x^\prime).$ Potom pro každé $x\in P_+(a,\delta)$ platí, že $a<x<x^\prime$, a tedy
	$$
	    m-\varepsilon < m \leq f(x) \leq f(x^\prime) < m+ \varepsilon.
	$$
	Potom máme $f(x)\in B(m,\varepsilon)$. Podobně postupujeme, pokud $m=-\infty$.
	\item Obdobně.\qedhere
\end{enumerate}
\end{proof}

\subsection{Funkce spojité na intervalu}
\begin{veta}[Bolzano]\label{veta3.10}
Nechť funkce $f$ je spojitá na intervalu $[a,b]$ a předpokládejme, že $f(a)<f(b)$. Potom pro každé $C\in (f(a),f(b))$ existuje $\xi\in(a,b)$ takové, že $f(\xi)=C$.
\end{veta}

\begin{proof}
	Nechť $C\in (f(a),f(b)).$ Položme $M = \left \{ z\in \left [ a,b \right ] , f(z)< C \right \} $.
	Množina $M$ je neprázdná ($a \in M$) a shora omezená, neboť $M \subseteq \left [ a,b \right ] $.
	Označme $\xi=\sup M$ (používáme vlastnost suprema). Přivedeme ke sporu možnost, že $f(\xi) < C$ i
	$f(\xi) < C$.
	\begin{enumerate}[($i$)]
		\item Pro spor předpokládejme, že $f(\xi)<C$. Platí tedy, že $\xi < b$. Funkce $f$ je v
		$\xi$ spojitá zprava. Podle věty \ref{veta3.5} (a) nalezneme $\delta \in \mathbb R ,\delta > 0$
		takové, že
		$$
		    \forall x \in \left [ \xi,\xi<\delta \right ):f(x)< C.
		$$
		Potom máme, že $\left [ \xi, \xi+ \delta \right ) \subset M $, což je spor.
		\item Pro spor předpokládejme, že $f(\xi) < C$. Potom $\xi > a$. Funkce $f$ je zleva spojitá
		v $\xi$. Podle věty \ref{veta3.5} (a) nalezneme $\delta \in \mathbb R , \delta > 0$
		takové, že
		$$\forall x \in (\xi-\delta, \xi]:f(x)>C.$$
		Potom $M\subseteq \left [ a,\xi-\delta \right ] $, což je spor ($\xi-\delta <\xi$  a $ \xi$ je
		horní závora).\qedhere
	\end{enumerate}
\end{proof}

\begin{poznamka}
	Věta \ref{veta3.10} platí i v případě, že $f(a) >f(b)$.
\end{poznamka}

\begin{lemma}\label{lemma3.11} % 3.11
Nechť $M\subset \mathbb R $ a platí
$$\forall x,y\in M\ \forall z\in \mathbb R :\ x<z<y\implies z\in M.$$
Pak $M$ je interval.
\end{lemma}

\begin{proof}
	Pokud $M = \emptyset$, pak $M$ je interval. Předpokládejme tedy, že $M\ne \emptyset$.
	Položme $a=\inf M, b =\sup M$. Ověříme, že interval $(a,b) \subseteq M \subseteq (a,b)\cup \left \{ a,b \right \} $.
	Nejprve ověříme $(a,b) \subseteq M$. Zvolme $z\in (a,b)$. Nalezneme prvky $x,y\in M$ splňující
	$x<z$ a $z<y$. Podle podmínky ze znění máme $z\in M$. Nyní ověříme $M \subseteq (a,b)\cup \left \{ a,b \right \} $.
	Zvolme $z\in M$. Potom $a\leq z\leq b$. Odtud plyne $z\in(a,b) \cup \left \{ a,b \right \} $.
\end{proof}

\begin{veta}\label{veta3.12} % 3.12
Nechť funkce $f\colon J\to\R$ je spojitá na intervalu $J$. Potom je $f(J)$ interval.
\end{veta}

\begin{proof}
	Použijeme lemma \ref{lemma3.11}. Předpokládejme, že $y_1,y_2 \in f(J)$ a $z\in \mathbb R $
	splňují $y_1 < z < y_2$. Nalezneme $x_1,x_2 \in J$ takové, že $f(x_1) = y_1, f(x_2) = y_2$.
	Podle věty \ref{veta3.10} nalezneme $\xi \in (x_1, x_2)$ (resp. $\xi \in (x_2, x_1)$) splňující
	$f(\xi) = z$ ($f$ je spojitá na $J$, a tedy i na $\left [ x_1,x_2 \right ] $, resp. $\left [ x_2,x_1 \right ] $).
	Máme tedy $z\in f(J)$. Podle lemmatu \ref{lemma3.11} je $f(J)$ interval.
\end{proof}

\begin{veta}\label{veta3.13}% 3.13
Nechť $f$ je spojitá funkce na intervalu $[a,b]$. Potom je $f$ na $[a,b]$ omezená.
\end{veta}

\begin{proof}
	Pro spor předpokládejme, že $f$ není omezená shora. Pro každé $n\in \mathbb N $
	nalezneme $x_n \in \left [ a,b \right ] $ takový, že $f(x_n) > n$. Díky Weierstrassově větě
	nalezneme vybranou posloupnost $\left \{ x_{n_k} \right \}_{k=1}^\infty $ z $\left \{ x_n \right \}_{n=1}^\infty $.
	Pak $\lim_{k \to \infty} x_{n_k} = x^* \in \left [ a,b \right ] $. Z varianty Heineovy věty plyne
	$\lim_{k \to \infty} f(x_{n_k})=f(x^*)\in \mathbb R $, $\forall k \in \mathbb N : f(x_{n_k}) > n_k \geq k$.
	odstud plyne $\lim_{k \to \infty} f(x_{n_k}) = \infty$, což je spor.
\end{proof}

\begin{poznamka}
	V důkazu věty \ref{veta3.13} jsme užili následující variantu Heineovy věty:
	Nechť $I\subset \mathbb R $ je interval a $f:I\to \mathbb R $. Pak je ekvivalentní:
	\begin{enumerate}[(i)]
		\item $f$ je spojitá na $I$,
		\item pro každé $c\in I$ a každou posloupnost $\left \{ x_n \right \} $ prvků intervalu $I$
		splňující $\lim_{n \to \infty} x_n=c$ platí $\lim_{n \to \infty} f(x_n)=f(c)$.
\end{enumerate}
\end{poznamka}

\begin{definice}
Nechť $M\subset\R$, $x\in M$ a funkce $f$ je definována alespoň na $M$ (tj. $M\subset\mathcal D(f)$).
Řekneme, že $f$ nabývá v bodě $x$ \textbf{maxima} (resp. \textbf{minima}) na $M$, jestliže platí
\begin{align*}
	\forall y\in M:\ f(y)\le f(x) && \textrm{(resp. }\forall y\in M:\ f(y)\ge f(x)\textrm{)}.
\end{align*}
Bod $x$ nazýváme \textbf{bodem maxima}
(resp. \textbf{minima}) funkce $f$ na množině $M$. Symbol $\max_M f$ (resp. $\min_M f$) označuje největší (resp. nejmenší) hodnotu,
které funkce $f$ na množině $M$ nabývá (pokud taková hodnota existuje).
\end{definice}

\begin{definice}
Nechť $M\subset \mathbb R $, $x\in M$ a funkce $f$ je definována alespoň na $M$ (tj. $M\subset \mathcal D(f)$).
Řekneme, že funkce $f$ má v bodě $x$
\begin{itemize}
	\item \textbf{lokální maximum vzhledem k $M$}, jestliže existuje $\delta>0$ takové,
	že pro každé $y\in P(x,\delta)\cap M$ platí $f(y)\le f(x)$,
	\item \textbf{lokální minimum vzhledem k $M$}, jestliže existuje $\delta>0$ takové,
	že pro každé $y\in P(x,\delta)\cap M$ platí $f(y)\geq f(x)$.
\end{itemize}
\end{definice}

\begin{veta}\label{veta3.14} % 3.14
Nechť $a,b\in \mathbb R $, $a<b$, a $f$ je spojitá funkce na intervalu $[a,b]$. Potom $f$ nabývá na $[a,b]$ svého maxima a svého minima.
\end{veta}

\begin{proof}
	Dokážeme existenci maxima. Podle věty \ref{veta3.13} je $f$ omezená na intervalu $\left [ a,b \right ] $,
	a proto $G=\sup f(\left [ a,b \right ] )$ je prvkem $\mathbb R $. Pro každé $n\in \mathbb N $
	nalezneme $x_n \in \left [ a,b \right ] $ splňující $f(x_n)> G-\frac{1}{n}$. Pomocí Weierstrassovy
	věty nalezneme vybranou posloupnost $\left \{ x_{n_k} \right \}_{k=1}^\infty $ z
	$\left \{ x_n \right \}_{n=1}^\infty $ takovou, že $\lim_{n \to \infty} x_{n_k}\in \mathbb R $.
	Potom platí, že $x^* = \lim x_{n_k} \in \left [ a,b \right ] $. Dále máme, že
	$\lim_{n \to \infty} f(x_n)=G$, neboť
	$$
	    \forall n \in \mathbb N:G-\frac{1}{n} < f(x_n) \leq G,
	$$
	a tedy $\lim_{n \to \infty} f(x_{n_k})=G$. Podle Heineovy věty máme, že $\lim_{n \to \infty} f(x_{n_k})=f(x^*)$.
	Odtud $f(x^*)=G$.
\end{proof}

\begin{veta}[O inverzní funkci]\label{veta3.15} % 3.15
Nechť $f$ je spojitá a rostoucí (resp. klesající) funkce na intervalu $J$. Potom funkce
$f^{-1}$ je spojitá a rostoucí (resp. klesjící) na intervalu $f(J)$.
\end{veta}

\begin{proof}
	Funkce $f$ je rostoucí, a tedy prostá. Funkce $f^{-1}$ je tedy dobře definovaná na $f(J)$.
	Množina $f(J)$ je interval podle věty \ref{veta3.12}. Je-li $f$ rostoucí (resp. klesající),
	pak je snadné ukázat, že $f^{-1}$ je rostoucí (resp. klesající). Předpokládejme, že $y_0\in f(J)$
	není pravým koncovým bodem intervalu $f(J)$. Budeme dokazovat, že $f^{-1}$ je spojitá zprava
	v bodě $y_0$. Označme $x_0= f^{-1}(y_0)$. Zvolme $\varepsilon \in \mathbb R , \varepsilon > 0$.
	Bod $y_0$ není pravý koncový bod $f(J)$, a tedy $x_0$ není pravý koncový bod $J$, neboť $f$
	je rostoucí. Nalezneme $x_1\in (x_0,x_0+\varepsilon) \cap J$. Nalezneme $\delta \in \mathbb R , \delta > 0$
	takové, že $f(x_1)=y_0+ \delta$. Pak pro každé $y\in \left [ y_0,y_0+\delta \right ) $
	platí $x_0=f^{-1}(y_0)\leq f^{-1}(y)\leq f^{-1}(y_0+\delta)=x_1$, takže
	$f^{-1}(y) \in [x_0,x_1) \subseteq B(x_0,\varepsilon)= B(f^{-1}(y_0),\varepsilon)$. Spojitost
	zleva v bodě, který není levým koncovým bodem $f(J)$, by se dokazovala analogicky.
\end{proof}
